\chapter{The Development Plan}
\label{ch:development-plan}

Chapter~\ref{ch:complete-pipeline} gives the logical order for layering features and models.
This chapter gives the actual build order: what to implement first given project priorities and foundation work that must happen before anything else.
Priority ordering: trading signal $>$ academic rigor $>$ model novelty.


%% ──────────────────────────────────────────────
\section{Milestones}
\label{sec:milestones}
%% ──────────────────────────────────────────────

Table~\ref{tab:milestones} defines eleven milestones with acceptance criteria and dependencies.
M0--M2 cover foundation work already underway; M3--M7 form the minimum viable deliverable; M8--M10 are upside.

\renewcommand{\arraystretch}{1.15}
\begin{table}[htbp]
\centering
\caption{Development milestones. M3--M7 form the minimum viable deliverable.}
\label{tab:milestones}
\footnotesize
\begin{tabular}{@{}p{2.5cm}p{5.2cm}p{3.5cm}c@{}}
\toprule
\textbf{Milestone} & \textbf{Acceptance Criteria} & \textbf{Key Tasks} & \textbf{Deps} \\
\midrule
\multicolumn{4}{@{}l@{}}{\small\itshape Foundation} \\[3pt]
\textbf{M0} Research \& Scoping &
  Literature survey (80+ papers); scope defined; feature taxonomy (L0--L5) designed; $\QLIKE$ chosen; LightGBM selected; universe (34 symbols) defined &
  Paper review, scope definition, metric/model selection &
  --- \\[6pt]
\textbf{M1} Data Infrastructure &
  Tick data via Chunk Store for 34 symbols; daily data from TSDB; IV surface from Marquee; resampling produces 78 bars/day; RV panel outputs 18 measures with parquet caching &
  Package scaffold, data integrations, resampling, RV panel builder, config system &
  M0 \\[6pt]
\textbf{M2} Feature Engine \& Baselines &
  Layer 0--1 features with no look-ahead; 7 HAR models fit and predict; $\QLIKE$ evaluation operational; purge gap $\geq h$ enforced; $\QLIKE$ sign matches \citet{Patton2011}; end-to-end pipeline runs; 390+ tests pass &
  HAR/asymmetry/noise-robust features, HAR family, metrics, CV splitters, CLI pipeline, purge gap fix, $\QLIKE$ sign fix &
  M1 \\
\midrule
\multicolumn{4}{@{}l@{}}{\small\itshape MVP} \\[3pt]
\textbf{M3} LightGBM &
  Custom $\QLIKE$ objective converges; Optuna finds improved params; walk-forward OOS predictions for 3 horizons &
  $\QLIKE$ gradient/hessian, model class, Optuna, walk-forward &
  M2 \\[6pt]
\textbf{M4} Tournament &
  8 models $\times$ 3 horizons table with DM $p$-values and Mincer--Zarnowitz efficiency tests on dev universe &
  Run baselines, DM test, MZ regression, tournament table &
  M3 \\[6pt]
\textbf{M5} Layer~2 Options &
  Options features produce daily values with no look-ahead; $\QLIKE$ lift documented &
  OptionsLayer, IV surface wiring, validation &
  M2 \\[6pt]
\textbf{M6} Layers 4--5 &
  Cross-asset and calendar features produce daily values; DY spillover index computed; FOMC/NFP/OPEX proximity indicators operational &
  Treasury slope, FX/commodity vol, DY spillover, event calendars &
  M2 \\[6pt]
\textbf{M7} Signal &
  IV--RV gap signal; equity curve; positive OOS Sharpe &
  Signal logic, P\&L backtest, performance metrics &
  M4, M5 \\
\midrule
\multicolumn{4}{@{}l@{}}{\small\itshape Upside} \\[3pt]
\textbf{M8} Ensemble &
  Residual stacking and prediction blending tested; 10-model tournament table &
  Residual stacking, inverse-$\QLIKE$ blending &
  M4 \\[6pt]
\textbf{M9} Microstructure \& Sequences &
  Layer~3 E-mini features (OBI, VPIN, depth); LSTM/TCN on intraday bars; scalar forecast as LightGBM feature evaluated &
  Microstructure layer, LSTM architecture, sequence pipeline &
  M2 \\[6pt]
\textbf{M10} Stretch &
  Ordered by impact: regime $\QLIKE$, MCS, Rashomon (TreeFARMS + RID + VIC), STreeD, Layer 6--7 (memory/sentiment), reporting/visualization, full universe, figures &
  Each task independent &
  M4--M9 \\
\bottomrule
\end{tabular}
\end{table}
\renewcommand{\arraystretch}{1.0}


%% ──────────────────────────────────────────────
\section{Critical Path}
\label{sec:critical-path}
%% ──────────────────────────────────────────────

Figure~\ref{fig:critical-path} shows milestone dependencies.
The critical path (M3 $\to$ M4 $\to$ M8 $\to$ M10) is highlighted.
M5, M6, and M9 branch from M2 and run in parallel with the critical path.
M7 requires both M4 (forecasts) and M5 (options features).

\begin{figure}[htbp]
\centering
\resizebox{\textwidth}{!}{%
\begin{tikzpicture}[
  >=Stealth,
  milestone/.style={draw, rounded corners, minimum height=0.9cm, minimum width=2.2cm,
    align=center, font=\small, fill=blue!8},
  critical/.style={milestone, line width=1.5pt, fill=orange!12},
  foundation/.style={milestone, dashed, fill=gray!8},
  arrow/.style={-{Stealth[length=2.5mm]}, thick},
  critarrow/.style={arrow, line width=1.5pt, color=defblue},
]
  % Foundation (dashed nodes)
  \node[foundation] (m0) at (0,0)   {M0\\Research};
  \node[foundation] (m1) at (3,0)   {M1\\Data};
  \node[foundation] (m2) at (6,0)   {M2\\Features};

  % Critical path (bold nodes)
  \node[critical] (m3) at (9,0)    {M3\\LightGBM};
  \node[critical] (m4) at (12,0)   {M4\\Tournament};
  \node[critical] (m8) at (15,0)   {M8\\Ensemble};
  \node[critical] (m10) at (18,0)  {M10\\Stretch};

  % Foundation arrows
  \draw[arrow] (m0) -- (m1);
  \draw[arrow] (m1) -- (m2);
  \draw[arrow] (m2) -- (m3);

  % Critical path arrows
  \draw[critarrow] (m3) -- (m4);
  \draw[critarrow] (m4) -- (m8);
  \draw[critarrow] (m8) -- (m10);

  % Parallel track (bottom row): M5, M6, M7, M9
  \node[milestone] (m5) at (9,-3)   {M5\\Options};
  \node[milestone] (m6) at (12,-3)  {M6\\Cross-Asset};
  \node[milestone] (m7) at (15,-3)  {M7\\Signal};
  \node[milestone] (m9) at (18,-3)  {M9\\Microstructure};

  % M2 fan-out: bus to M5, M6, M9 (all depend on M2)
  \draw (m2.south) -- ++(0,-1.05) coordinate (branch);
  \fill (branch) circle (2pt);
  \draw (branch) -- (branch -| m9.north);
  \draw[arrow] (branch -| m5.north) -- (m5.north);
  \draw[arrow] (branch -| m6.north) -- (m6.north);
  \draw[arrow] (branch -| m9.north) -- (m9.north);

  % M4 -> M7 (M7 depends on M4): route above bus then down
  \draw[arrow] (m4.south) -- ++(0,-0.65) -| (m7.north);

  % M5 -> M7 (M7 depends on M5): route above M6
  \draw[arrow] (m5.north east) -- ++(0,0.5) -| (m7.north west);

  % M9 -> M10 (M10 depends on M4--M9): clean vertical
  \draw[arrow] ([xshift=4mm]m9.north) -- ([xshift=4mm]m10.south);

  % M6 -> M10 (M10 depends on M4--M9): route below bottom row, up through M7/M9 gap
  \draw[arrow] (m6.south) -- ++(0,-0.5) -- (16.5,-3.95) -- (16.5,-0.45) -- (m10.south);
\end{tikzpicture}%
}
\caption{Milestone dependencies.
Dashed nodes (M0--M2): foundation work.
Bold path: critical path (M3--M4--M8--M10).
M5, M6, and M9 branch from M2 and run in parallel.
M7 requires both M4 (forecasts) and M5 (options features).
M10 waits on all prior milestones.}
\label{fig:critical-path}
\end{figure}


%% ──────────────────────────────────────────────
%% Boxes
%% ──────────────────────────────────────────────

\begin{keyidea}[M3--M7 Is the Minimum Viable Deliverable]
A $\QLIKE$ tournament (7~HAR variants + LightGBM, DM tests) plus cross-asset features and a tradeable IV--RV signal with P\&L backtest is a presentable result.
Everything in M8--M10 is upside.
If time runs out after M7, the project has a defensible outcome.
\end{keyidea}

\begin{warning}[M2 Must Be Complete Before M3]
The purge gap bug causes silent data leakage for $h=22$.
The $\QLIKE$ sign convention determines whether the loss function penalizes over-prediction or under-prediction correctly.
All results produced before M2 is complete are potentially invalid.
Do not skip ahead.
\end{warning}
